\begin{abstract}
We identify a gradient-scale failure mode of per-neuron Lagrangian
FLOPs pruning in Mixture-of-Experts (MoE) vision transformers: the
per-neuron penalty gradients are orders of magnitude weaker than task
gradients (${\sim}10^{-4}$ vs.\ ${\sim}10^{-2}$), causing the
optimizer to satisfy the budget constraint through \emph{uniform mask
collapse} rather than selective pruning. We validate this prediction
empirically: a properly tuned Lagrangian baseline achieves only
$10.4\%$ accuracy on CIFAR-100 (barely above the $1\%$ random
baseline), $36.4\%$ on CIFAR-10, and $4.7\%$ on Tiny-ImageNet---despite
reaching the target FLOPs ratio on every dataset. We further show that
\emph{any} schedule-based pruning method (Linear, Random, Magnitude,
or our cubic variant) bypasses this failure mode and recovers
near-dense accuracy, indicating that the diagnosis---not a particular
schedule---is the primary contribution. As a practical remedy we
propose \textbf{KaleidoNet}, which replaces the learned Lagrangian
penalty with a deterministic \emph{cubic sparsity schedule}, gradient
masking, and dual-rate optimization. Across all three datasets trained
to convergence (50{,}000 steps, 3 seeds), KaleidoNet retains
98.3--99.2\% of dense-model accuracy with $1.8{\times}$ FLOPs
reduction and $1.63{\times}$ parameter compression, the strongest
schedule-based instantiation we tested. Ablations over eight
configurations isolate the contribution of each component, and model
surgery converts soft masks into physically compact networks. We
report a current implementation limitation: at this scale MoE routing
overhead dominates wall-clock latency on Apple MPS, so the FLOPs and
parameter savings do not yet translate into a deployment speedup; we
discuss this as a platform-specific limitation rather than a property
of the algorithm. Validation is limited to small-to-medium scale
(CIFAR-10/100, Tiny-ImageNet; 5.49M-parameter MoE ViT); ImageNet-1K
validation with larger backbones remains future work.
\end{abstract}
